% Clase 3 — Razonamiento bajo Incertidumbre (parte 1)
% Lectura: AIMA Cap. 12
\section{Razonamiento bajo Incertidumbre I}

\subsection{Probabilidad como Base del Razonamiento}

Cuando el entorno es \textbf{parcialmente observable} o \textbf{estocástico}, se usa probabilidad para representar la incertidumbre.

\begin{conceptbox}
La \textbf{probabilidad a priori} (o incondicional) $P(A)$ representa el grado de creencia sin ninguna evidencia adicional.

La \textbf{probabilidad condicional} $P(A \mid B) = \dfrac{P(A \land B)}{P(B)}$ representa el grado de creencia dado que $B$ es verdadero.
\end{conceptbox}

\subsection{Regla de Bayes}

$$P(A \mid B) = \frac{P(B \mid A) \cdot P(A)}{P(B)}$$

\begin{itemize}
    \item $P(A)$: probabilidad a priori de $A$ (hipótesis)
    \item $P(B \mid A)$: verosimilitud de la evidencia $B$ dado $A$
    \item $P(B)$: probabilidad marginal de $B$ (normalizador)
    \item $P(A \mid B)$: probabilidad a posteriori
\end{itemize}

\subsection{Distribuciones de Probabilidad}

Para variables aleatorias discretas, la \textbf{distribución conjunta completa} especifica $P(X_1, X_2, \ldots, X_n)$ para todos los valores posibles. De ella se puede derivar cualquier probabilidad por marginalización y condicionamiento.

\textbf{Marginalización}: $P(X) = \sum_y P(X, Y=y)$

\textbf{Independencia}: $X \perp Y \Leftrightarrow P(X, Y) = P(X) \cdot P(Y)$

\textbf{Independencia condicional}: $X \perp Y \mid Z \Leftrightarrow P(X \mid Y, Z) = P(X \mid Z)$
